\documentclass{article}

\usepackage{amsmath,amssymb}
\usepackage{xurl}
\usepackage[hidelinks]{hyperref}
\setlength{\emergencystretch}{2em}

\input{command}

\title{The Rectangular Choi--Jamiolkowski Correspondence:\\
Resolution of the Square-Dimensional Scope Restriction}
\author{}
\date{2026-08-04}

\begin{document}
\maketitle
\gapnote{scope-restriction}{resolved}

\section{The source statement}

Let $T:\MN{d}\to\MN{d'}$ be linear. Let
$\Omega_d\in\C^d\otimes\C^d$ be the normalized maximally entangled vector,
and let $\tau_T$ denote the normalized Choi matrix. Proposition~2.1 of
Ref.~\cite{Wolf2012Quantum} defines
\begin{align}
    \tau_T
        &=(T\otimes\operatorname{id}_d)
          (\ketbra{\Omega_d}{\Omega_d})
          \in \MN{d'}\otimes\MN{d}.
    \label{eq:wolf_choi_rectangular_scope_choi_matrix}
\end{align}
For $A\in\MN{d'}$ and $B\in\MN{d}$, it gives the trace pairing
\begin{align}
    \tr(A\,T(B))
        &=d\,\tr(\tau_T(A\otimes B^T)).
    \label{eq:wolf_choi_rectangular_scope_trace_pairing}
\end{align}
These formulas define mutually inverse correspondences between linear maps
and bipartite operators \cite[Proposition~2.1]{Wolf2012Quantum}.

The same proposition identifies Hermiticity preservation with
\begin{align}
    \tau_T
        &=\tau_T^\dagger,
    \label{eq:wolf_choi_rectangular_scope_hermiticity_condition}
\end{align}
and complete positivity with
\begin{align}
    \tau_T
        &\geq0.
    \label{eq:wolf_choi_rectangular_scope_complete_positivity_condition}
\end{align}
Writing $\tr_A$ and $\tr_B$ for the partial traces over the output and input
factors, respectively, unitality is equivalent to
\begin{align}
    \tr_B(\tau_T)
        &=\frac{\Id_{d'}}{d},
    \label{eq:wolf_choi_rectangular_scope_unital_partial_trace}
\end{align}
and trace preservation is equivalent to
\begin{align}
    \tr_A(\tau_T)
        &=\frac{\Id_d}{d}.
    \label{eq:wolf_choi_rectangular_scope_trace_preserving_partial_trace}
\end{align}
Without assuming trace preservation, Proposition~2.1 of
Ref.~\cite{Wolf2012Quantum} gives
\begin{align}
    \tr_A(\tau_T)
        &=\frac{(T^*(\Id_{d'}))^T}{d},
    \label{eq:wolf_choi_rectangular_scope_general_left_partial_trace}\\
    \tr(\tau_T)
        &=\frac{\tr(T^*(\Id_{d'}))}{d}.
    \label{eq:wolf_choi_rectangular_scope_general_trace}
\end{align}
It also gives the corresponding proportional-identity criterion for a
doubly stochastic map \cite[Proposition~2.1]{Wolf2012Quantum}.

Theorem~2.1 of Ref.~\cite{Wolf2012Quantum} states that $T$ is completely
positive if and only if it has a rectangular Kraus representation
\begin{align}
    T(X)
        &=\sum_{j=1}^{r}K_jXK_j^\dagger,
    \label{eq:wolf_choi_rectangular_scope_rectangular_kraus_representation}\\
    K_j
        &\in M_{d'\times d}(\C).
    \label{eq:wolf_choi_rectangular_scope_rectangular_kraus_dimensions}
\end{align}
The theorem also identifies the trace-preserving and unital normalization
conditions, the minimal Kraus cardinality
\begin{align}
    r
        &=\operatorname{rank}(\tau_T)
        \leq dd',
    \label{eq:wolf_choi_rectangular_scope_minimal_kraus_rank}
\end{align}
an orthogonal minimal family, and unitary freedom after padding the smaller
family with zero operators \cite[Theorem~2.1]{Wolf2012Quantum}.

\section{Resolution}

The different-dimension Choi statements are formalized in the namespace
\leanid{ChoiRectangular}.\footnote{Formal module:
    \doclink{QICLean/Channel/ChoiRectangular}.}
The development defines \leanid{choiMatrix} on
$\C^{d'}\otimes\C^d$ and proves:
\begin{itemize}
    \item the inverse identities \leanid{mapOfChoiMatrix_choiMatrix} and
        \leanid{choiMatrix_mapOfChoiMatrix};

    \item the trace-pairing formula \leanid{trace_pairing};

    \item the Hermiticity equivalence
        \leanid{choiMatrix_isHermitian_iff_hermiticityPreserving};

    \item the complete-positivity equivalence
        \leanid{isKrausCP_iff_choiMatrix_posSemidef};

    \item the unitality equivalence
        \leanid{unital_iff_traceRight_choiMatrix}, obtained from
        \leanid{traceRight_choiMatrix} and
        (\ref{eq:wolf_choi_rectangular_scope_unital_partial_trace});

    \item the trace-preservation equivalence
        \leanid{tracePreserving_iff_traceLeft_choiMatrix}, obtained from
        \leanid{traceLeft_choiMatrix} and
        (\ref{eq:wolf_choi_rectangular_scope_general_left_partial_trace});

    \item the normalization formula \leanid{trace_choiMatrix}, corresponding
        to (\ref{eq:wolf_choi_rectangular_scope_general_trace});

    \item the doubly-stochastic equivalence
        \leanid{doublyStochastic_iff_partialTraces_proportional}.
\end{itemize}

The rectangular Kraus development is in
\doclink{QICLean/Channel/KrausRectangular}, in the same namespace. It proves
the normalization equivalences
\leanid{kraus_tp_iff_sum_conjTranspose_mul} and
\leanid{kraus_unital_iff_sum_mul_conjTranspose}. It also proves minimality of
the Kraus rank through
\leanid{choiRank_isLeast_hasKrausCard_of_isKrausCP}, the bound
\leanid{choiRank_le_mul}, and existence of an orthogonal minimal family
through \leanid{exists_kraus_orthogonal_of_isKrausCP}. The padded unitary
freedom is formalized by the dimension-generic declarations
\leanid{kraus_isometry_freedom_iff} and
\leanid{kraus_unitary_freedom_iff} in
\doclink{QICLean/Channel/KrausUnitaryFreedom}.

The complete-positivity predicate is \leanid{IsKrausCP}: the existence of
rectangular Kraus operators. This is the complete-positivity notion used
elsewhere in the development, and the Choi equivalence is stated relative to
it.

\section{Surviving square-only declarations}

The dimension-generic Kraus API states cardinality, rank, and freedom results
for rectangular operators $K_j\in M_{d'\times d}(\C)$. The square development
is the specialization $d=d'$.

The declarations \leanid{Channel.HasKrausCard},
\leanid{Channel.HasKrausRankLE}, \leanid{Channel.choiRank},
\leanid{Channel.hasKrausCard_mono},
\leanid{Channel.choiMatrix_eq_sum_vecMulVec_of_kraus},
\leanid{Channel.choiRank_le_of_hasKrausCard}, and
\leanid{Channel.hasKrausCard_choiRank_of_cp} are in
\doclink{QICLean/Channel/KrausRank}. The normalization declarations
\leanid{kraus_sum_conjTranspose_mul_of_tp},
\leanid{kraus_sum_mul_conjTranspose_of_unital}, and
\leanid{kraus_same_map_of_isometry_combination} are in
\doclink{QICLean/Channel/KrausRepresentation}. The dual, Gramian, and
necessary-direction freedom declarations
\leanid{kraus_dual_eq_of_map_eq},
\leanid{kraus_conjTranspose_mul_eq_of_map_eq},
\leanid{kraus_rectangular_freedom}, and
\leanid{kraus_rectangular_freedom'} are in
\doclink{QICLean/Channel/KrausFreedom}. The freedom equivalences
\leanid{kraus_isometry_freedom_iff} and
\leanid{kraus_unitary_freedom_iff} are in
\doclink{QICLean/Channel/KrausUnitaryFreedom}. Their former duplicates in the
\leanid{ChoiRectangular} namespace have been removed.

The following declarations remain square-only. They are specializations of
rectangular or dimension-generic results, not independent scope gaps.
\begin{itemize}
    \item \leanid{ChoiJamiolkowski.choiMatrix} and its correspondence clauses
        remain square. The square Choi matrix is definitionally the
        rectangular matrix at $d=d'$, as stated by
        \leanid{ChoiRectangular.choiMatrix_eq_choiJamiolkowski}. Each square
        clause specializes the rectangular clause with the same name.

    \item \leanid{IsTracePreservingMap} and \leanid{IsChannel} are stated for
        endomorphisms of one matrix algebra. The rectangular development
        expresses trace preservation pointwise as
        \begin{align}
            \tr(T(X))
                &=\tr(X),
            \label{eq:wolf_choi_rectangular_scope_pointwise_trace_preservation}
        \end{align}
        and represents rectangular complete positivity together with trace
        preservation by \leanid{IsKrausCPTP}.

    \item \leanid{Channel.hasKrausRankLE_choiRank_of_cptp} and
        \leanid{kraus_transition_unitary_of_hs_orthonormal} have no
        rectangular counterparts. The first uses the square-only
        \leanid{IsChannel} structure. The second is a square linear-algebraic
        core lemma with square-only call sites.
\end{itemize}

\section{Conclusion}

The scope restriction recorded on 2026-07-30 is lifted. Proposition~2.1 and
Theorem~2.1 of Ref.~\cite{Wolf2012Quantum} are formalized with independent
input and output dimensions and with all displayed clauses. No strengthened
hypothesis or source-level gap remains in this area. The resolution is tracked
in
\href{https://github.com/LionSR/TNLean/issues/3987}{TNLean issue \#3987}.

\bibliographystyle{alpha}
\bibliography{references}

\end{document}
